\documentclass[a4paper, 12pt]{article}

%====================== PACKAGES ======================

\usepackage[top=2cm, bottom=2cm, left=2cm, right=2cm]{geometry}
\usepackage[french]{babel}
\usepackage[utf8x]{inputenc}
\usepackage{titlesec}
\usepackage{titletoc,tocloft}
\usepackage{float}
\usepackage[export]{adjustbox}
\usepackage{amsmath}
\usepackage{wrapfig}
\usepackage{graphicx}
\usepackage{url}
\usepackage{array}
\usepackage{tabularx}
\usepackage{setspace}
\usepackage{abstract}
\usepackage[T1]{fontenc}
\usepackage{xcolor}
\usepackage{subfig}
\usepackage{moresize}
\usepackage{hyperref}
\usepackage{minted}

%====================== INFORMATION ET REGLES ======================

\hypersetup{							% Information sur le document
pdfauthor = {Aurélien Verbeke},			% Auteurs
pdftitle = {Rapport projet minishell SEC}			% Titre du document
}

%code
\definecolor{couleur_fond_code_fonce}{HTML}{272822}
\definecolor{couleur_fond_code_clair}{HTML}{EAEAEA}
\newminted{c}{autogobble, samepage, bgcolor=couleur_fond_code_clair, style=staroffice}
\newmintedfile{c}{autogobble, samepage, bgcolor=couleur_fond_code_fonce, style=monokai}
\newmintinline{c}{style=staroffice, bgcolor=couleur_fond_code_clair}
\newminted{bash}{autogobble, samepage, bgcolor=couleur_fond_code_clair, style=staroffice, fontsize=\tiny}

% a chaque nouvelle partie, la numérotation des sections repart de zéro
\makeatletter
\@addtoreset{section}{part}
\makeatother

% affichage des titres
\titlespacing*{\section}{0cm}{2.5cm}{0.5cm}
\titlespacing*{\part}{0cm}{1cm}{1cm}
\titleformat{\part}[hang]{\bfseries\Huge}{\thepart \ --}{12pt}{}
\titleformat*{\section}{\Large\bfseries}

\setcounter{tocdepth}{2} % dans la table des matières, on s'arrête aux sous-sections

% affichage des différentes tables
\setlength{\cftsecindent}{0.5cm}
\setlength{\cftsubsecindent}{1cm}
\renewcommand{\cftbeforesecskip}{0.3cm}
\renewcommand{\cftfigpresnum}{\textsc{Figure} \ }
\setlength{\cftfignumwidth}{2.5cm}
\renewcommand{\cfttabpresnum}{\textsc{Table} \ }
\setlength{\cfttabnumwidth}{2.3cm}

%espacement entre les lignes d'un tableau
\renewcommand{\arraystretch}{1.5}

%régler l'espacement entre les lignes
\newcommand{\HRule}{\rule{\linewidth}{1.5mm}}

%======================== DEBUT DU DOCUMENT ========================
\begin{document}

%page de garde
\renewcommand\thepage{} % ne pas numéroter

\begin{titlepage}
\begin{center}

% Upper part of the page. The '~' is needed because only works if a paragraph has started.
\textsc{\Huge Enseeiht}\\[1cm]

\vfill

% Title
\HRule

\textsc{\huge{\textbf{Rapport projet SEC\\[1cm]Minishell}}}

\HRule

\vfill

% Author and supervisor
\Large{Aurélien \textsc{Verbeke}}

\vspace{1cm}

\large{Année 2024 - 2025}

\vspace{2cm}

\includegraphics[height=2.cm, valign=m]{3-Toulouse-INP-N7-couleur-sans-fond.png}

\vspace*{-0.5cm}

\end{center}
\end{titlepage}

%sommaire
\setcounter{page}{1}
\renewcommand\thepage{\arabic{page}}
\tableofcontents

%====================== INCLUSION DES PARTIES ======================

\newpage
\part{Introduction}

Le \textsc{Minishell} est un shell minimaliste qui permet de lancer des commandes, changer de répertoire courant, afficher le contenu d'un répertoire.

Les commandes peuvent être lancées en arrière-plan et chaînées, et l'utilisateur peut faire des redirections d'entrées / sorties.

\part{Processus}

\noindent
\textbf{Etape 2}

\bigskip
On ajoute \cinline/execvp(commande[0], commande)/. Cela permet de remplacer le code courant par celui de la commande.

On utilise la version \texttt{vp} pour pouvoir transmettre tous les arguments en une seule fois et pour pouvoir utiliser le \texttt{PATH}.

\bigskip
\noindent
\textbf{Etape 3}

\bigskip
Dans le père, on ajoute la directive \cinline/wait()/ pour attendre la termiaison du fils lancé.

\bigskip
\noindent
\textbf{Etape 4}

\bigskip
Dans le père, on remplace \cinline/wait()/ par
\begin{ccode}
if (ligneCommande->backgrounded == NULL) {
	int status;
	if (waitpid(pidFils, &status, 0) == -1) {
		exit(EXIT_FAILURE);
	}
}
\end{ccode}

On attendra alors que l'éxécution des fils en avant-plan se termine.

Il est nécessaire d'effectuer un \cinline/waitpid()/ sur le fils lancé pour ne pas que le traitement du père soit rétabli après la terminaison du fils en arrière-plan.

\newpage
\part{Signaux}

\noindent
\textbf{Question}

\bigskip
Voir étape 4.

\bigskip
\noindent
\textbf{Etape 5}

\bigskip
On ajoute en début de \cinline/main()/ :

\begin{ccode}
struct sigaction actionFilsTermine;
actionFilsTermine.sa_handler = handlerFilsTermine;
actionFilsTermine.sa_flags = SA_RESTART;
sigemptyset(&actionFilsTermine.sa_mask);
if (sigaction(SIGCHLD, &actionFilsTermine, NULL) == -1) {
	exit(EXIT_FAILURE);
}
\end{ccode}

La fonction de traitement est

\begin{ccode}
void handlerFilsTermine() {
	fprintf(stderr, "Handler : Le fils a change d'etat\n");
}
\end{ccode}

\bigskip
\noindent
\textbf{Etape 6}

\bigskip
On peut remplacer la fonction de traitement par
\begin{ccode}
void handlerFilsTermine() {
	int status;
	pid_t pid = waitpid(-1, &status, WNOHANG|WUNTRACED|WCONTINUED);

	fprintf(stderr, "Handler : Fils %d termine\n", pid);
}
\end{ccode}

\bigskip
\noindent
\textbf{Etape 7}

\bigskip
On remplace le \cinline/waitpid()/ du traitement du père par un simple \cinline/pause()/.

On constate que le terminal rend la main après le \texttt{sleep 10 &}. C'est normal puisque on attend n'importe quel signal, en particulier un pouvant venir d'un fils en arrière plan.

\bigskip
\noindent
\textbf{Etape 8}

\bigskip
L'envoi du signal \texttt{SIGSTOP} fait entrer le processus dans un état de pause (\texttt{T+}).

L'envoi du signal \texttt{SIGCONT} relance le processus (état \texttt{S+} puisque \texttt{sleep}).

\bigskip
\noindent
\textbf{Etape 9}

\bigskip
On remplace le simple affichage de la fonction de traitement \cinline/handlerFilsTermine()/ par :

\begin{ccode}
if (WIFEXITED(status)) {
	fprintf(stderr, "Handler : Fils %d termine avec le code %d\n",
					pid, WEXITSTATUS(status));
} else if (WIFSTOPPED(status)) {
	fprintf(stderr, "Handler : Fils %d suspendu\n", pid);
} else if (WIFCONTINUED(status)) {
	fprintf(stderr, "Handler : Fils %d repris\n", pid);
} else if (WIFSIGNALED(status)) {
	fprintf(stderr, "Handler : Fils %d tue\n", pid);
}
\end{ccode}

\noindent
\textbf{Résultats}

Pour le lancement de \texttt{sleep 60 &} et les envois consécutifs des signaux \texttt{SIGSTOP} et \texttt{SIGCONT} :

\begin{bashcode}
> Handler : Fils 48136 suspendu
Handler : Fils 48136 repris
Handler : Fils 48136 termine avec le code 0
\end{bashcode}

\begin{bashcode}
ave9761@forge:~/git_perso/1a/sec/rapport$ ps -fjn
   PPID     PID    PGID     SID TTY        TPGID STAT   UID   TIME COMMAND
  41282   41289   41289   41289 pts/1      48123 Ss   57562   0:00 bash -c source ~/.bashrc; $SHELL -i
  41289   41292   41292   41289 pts/1      48123 S    57562   0:00  \_ /bin/bash -i
  41292   48123   48123   41289 pts/1      48123 S+   57562   0:00      \_ ./minishell
  48123   48136   48123   41289 pts/1      48123 S+   57562   0:00          \_ sleep 20
  40174   40176   40176   40176 pts/0      48146 Ss   57562   0:00 bash -c source ~/.bashrc; $SHELL -i
  40176   40179   40179   40176 pts/0      48146 S    57562   0:00  \_ /bin/bash -i
  40179   48146   48146   40176 pts/0      48146 R+   57562   0:00      \_ ps -fjn
ave9761@forge:~/git_perso/1a/sec/rapport$ kill -STOP 48136
ave9761@forge:~/git_perso/1a/sec/rapport$ ps -fjn
   PPID     PID    PGID     SID TTY        TPGID STAT   UID   TIME COMMAND
  41282   41289   41289   41289 pts/1      48123 Ss   57562   0:00 bash -c source ~/.bashrc; $SHELL -i
  41289   41292   41292   41289 pts/1      48123 S    57562   0:00  \_ /bin/bash -i
  41292   48123   48123   41289 pts/1      48123 S+   57562   0:00      \_ ./minishell
  48123   48136   48123   41289 pts/1      48123 T+   57562   0:00          \_ sleep 20
  40174   40176   40176   40176 pts/0      48167 Ss   57562   0:00 bash -c source ~/.bashrc; $SHELL -i
  40176   40179   40179   40176 pts/0      48167 S    57562   0:00  \_ /bin/bash -i
  40179   48167   48167   40176 pts/0      48167 R+   57562   0:00      \_ ps -fjn
ave9761@forge:~/git_perso/1a/sec/rapport$ kill -CONT 48136
ave9761@forge:~/git_perso/1a/sec/rapport$ ps -fjn
   PPID     PID    PGID     SID TTY        TPGID STAT   UID   TIME COMMAND
  41282   41289   41289   41289 pts/1      48123 Ss   57562   0:00 bash -c source ~/.bashrc; $SHELL -i
  41289   41292   41292   41289 pts/1      48123 S    57562   0:00  \_ /bin/bash -i
  41292   48123   48123   41289 pts/1      48123 S+   57562   0:00      \_ ./minishell
  48123   48136   48123   41289 pts/1      48123 S+   57562   0:00          \_ sleep 20
  40174   40176   40176   40176 pts/0      48188 Ss   57562   0:00 bash -c source ~/.bashrc; $SHELL -i
  40176   40179   40179   40176 pts/0      48188 S    57562   0:00  \_ /bin/bash -i
  40179   48188   48188   40176 pts/0      48188 R+   57562   0:00      \_ ps -fjn
\end{bashcode}

\bigskip
\noindent
\textbf{Etape 10}

\bigskip
Que le processus soit en arrière-plan ou en avant-plan, on tue le minishell et la commande lancée.

\newpage
\noindent
\textbf{Etape 11.1}

\bigskip
On ajoute :

\begin{ccode}
struct sigaction actionCtrlCZ;
actionCtrlCZ.sa_handler = handlerCtrlCZ;
actionCtrlCZ.sa_flags = SA_RESTART;
sigemptyset(&actionCtrlCZ.sa_mask);

if (sigaction(SIGTSTP, &actionCtrlCZ, NULL) == -1) {
	exit(EXIT_FAILURE);
}
if (sigaction(SIGINT, &actionCtrlCZ, NULL) == -1) {
	exit(EXIT_FAILURE);
}
\end{ccode}

La fonction de traitement est :

\begin{ccode}
void handlerCtrlCZ() {
	fprintf(stderr, "Signal ignore par le pere.");
}
\end{ccode}

Le minishell bloque alors tout \texttt{Ctrl-C} et \texttt{Ctrl-Z} pour les fils et lui-même, ce qui est problématique.

\bigskip
\noindent
\textbf{Etape 11.2}

\bigskip
De même qu'à l'étape précédente, on associe le traitement \texttt{SIG\_IGN} aux signaux \texttt{SIGTSTP} et \texttt{SIGINT} avant le \cinline/fork()/.

Il faut alors, dans le fils, réassocier les traitements initiaux avec \texttt{SIG\_DFL} aux signaux \texttt{SIGTSTP} et \texttt{SIGINT} pour qu'ils continue à les recevoir.

Pour plus de cohérence, on continue d'utiliser la fonction \cinline/sigaction()/.

\bigskip
Le résultat obtenu est celui qui était attendu : seulement les processus fils sont suspendus ou stoppés.

\bigskip
\noindent
\textbf{Etape 11.3}

\bigskip
On a :

\begin{ccode}
sigset_t masquePere;
sigset_t masqueFils;

sigemptyset(&masqueFils);

sigemptyset(&masquePere);
sigaddset(&masquePere, SIGINT);
sigaddset(&masquePere, SIGTSTP);
\end{ccode}

Et on ajoute \cinline/sigprocmask(SIG_SETMASK, &masquePere, NULL);/ avant le \cinline/fork()/ et \cinline/sigprocmask(SIG_SETMASK, &masqueFils, NULL);/ dans le fils.

Le comportement est alors le même qu'à l'étape précédente, comme espéré.

\newpage
\noindent
\textbf{Etape 12}

On effectue \cinline/setpgrp()/ dans le fils uniquement dans le cas où la commande est lancée en arrière-plan. Le fils est alors détaché du groupe de son père et ne reçoit plus les signaux correspondant à \texttt{Ctrl-C} et \texttt{Ctrl-Z}.

\part{Fichiers et redirections}

\bigskip
\noindent
\textbf{Etape 13}

\bigskip
Pour la redirection en entrée : il faut ouvrir le fichier donné\\
\noindent(\cinline/fdSource = open(ligneCommande->in, O_RDONLY)/) et remplacer le descripteur de l'entrée standard par celui du fichier ouvert (\cinline/dup2(fdSource, 0)/).

Pour la redirection en sortie : il faut ouvrir le fichier donné et le créer s'il n'existe pas (\cinline/fdDest = open(ligneCommande->out, O_WRONLY|O_CREAT|O_TRUNC, 0644)/) et remplacer le descripteur de la sortie standard par celui du fichier ouvert (\cinline/dup2(fdDest, 1)/).

\bigskip
\noindent
\textbf{Etape 14}

\bigskip
On crée d'abord une fonction qui change de répertoire :

\begin{ccode}
void changerRepertoireCourant(const char* _nouveauRepertoire) {
	char erreur[512];
	char* nouveauRepertoire;

	// pas de repertoire fourni, on va dans le HOME
	if (_nouveauRepertoire == NULL) {
		nouveauRepertoire = getenv("HOME");
		sprintf(erreur, "Erreur changement repertoire vers home\n");
	}

	// on va dans le repertoire fourni
	else {
		nouveauRepertoire = malloc(256*sizeof(char));
		strcpy(nouveauRepertoire, _nouveauRepertoire);
		sprintf(erreur, "Erreur changement repertoire vers %s\n",
							nouveauRepertoire);
	}

	// execution du changement de repertoire
	if (chdir(nouveauRepertoire) == -1) {
		fprintf(stderr, erreur);
		exit(EXIT_FAILURE);
	}
}
\end{ccode}

\newpage
On ajoute ensuite la prise en compte de la commande dans la boucle principale :

\begin{ccode}
else if (strcmp(commande[0], "cd") == 0) {
	changerRepertoireCourant(commande[1]);
}
\end{ccode}

Enfin, on modifie l'affichage du prompt pour savoir où on est :

\begin{ccode}
printf("(%s) > ", getcwd(NULL, 0));
\end{ccode}

\bigskip
\noindent
\textbf{Etape 15}

\bigskip
De même, on crée d'abord une fonction qui affiche un répertoire :

\begin{ccode}
void afficherRepertoire(const char* _repertoire) {
        char erreur[512];
        char* repertoire;
        DIR* entreesRepertoire;
        struct dirent* entreeRepertoire;

        // pas de repertoire fourni, on affiche le repertoire courant
	if (_repertoire == NULL) {
                repertoire = getcwd(NULL, 0);
                sprintf(erreur, "Erreur ouverture repertoire courant\n");
        }

	// on affiche le repertoire fourni
        else {
                repertoire = malloc(256*sizeof(char));
                strcpy(repertoire, _repertoire);
                sprintf(erreur, "Erreur ouverture repertoire %s\n", repertoire);
        }

	// ouverture du repertoire
        if ((entreesRepertoire = opendir(repertoire)) == NULL) {
                fprintf(stderr, erreur);
                exit(EXIT_FAILURE);
        }

	// affichage du contenu
        while ((entreeRepertoire = readdir(entreesRepertoire)) != NULL) {
                printf("%s\n", entreeRepertoire->d_name);
        }

	// fermeture du repertoire
        if (closedir(entreesRepertoire) == -1) {
                fprintf(stderr, "Erreur fermeture repertoire\n");
                exit(EXIT_FAILURE);
        };
}
\end{ccode}

\newpage
Puis on ajoute la prise en compte de la commande dans la boucle principale :

\begin{ccode}
else if (strcmp(commande[0], "dir") == 0) {
	afficherRepertoire(commande[1]);
}
\end{ccode}

\part{Tubes}

\noindent
\textbf{Etapes 16 et 17}

\bigskip
Avant la création du fils pour l'éxécution de la commande, si une commande pipelinée suit :
\begin{itemize}
	\item On sauvegarde le tube précédent
	\item On crée un nouveau tube
\end{itemize}

Puis, dans le fils :
\begin{itemize}
	\item Une commande pipelinée suit : on ferme le tube en lecture, on remplace la sortie standard par le tube puis on ferme le tube en écriture.
	\item La commande est au moins la deuxième du pipeline : on ferme l'ancien tube en écriture, on remplace l'entrée standard par l'ancien tube puis on ferme l'ancien tube en lecture.
\end{itemize}

Cependant, il faut gérer les tubes dans le père. En effet, il ne faut pas que le père reste rédacteur pour le tube. Sinon, la lecture est bloquante. Il doit cependant rester lecteur pour que la commande suivante puisse remplacer son entrée standard par la sortie du tube.

\bigskip
\textbf{Remarque : }
La conception ainsi réalisée ne change donc pas le mécanisme de lancement des commandes.

\bigskip
Résultat pour 2 terminaux ouverts :
\begin{bashcode}
(/home/ave9761/git_perso/1a/sec/minishell) > who | grep ave9761 | wc -l
Handler : Fils 220861 termine avec le code 0
Handler : Fils 220862 termine avec le code 0
2
Handler : Fils 220863 termine avec le code 0
(/home/ave9761/git_perso/1a/sec/minishell) >
\end{bashcode}

\part{Remarques supplémentaires}

On fait bien attention à ne pas faire de \cinline/printf()/ pour ce qui ne concerne pas les commandes (informations, erreurs...).
Sinon, on risque d'injecter des données dans les canaux \cinline/stdin/ (0) et \cinline/stdout/ (1).

\bigskip
Dans le fils, il faut ajouter toutes les commandes avant le \texttt{exec}. Sinon, elles ne sont jamais éxécutées et remplacées par le code de la commande lancée.

\end{document}
